Guided-wave structural health monitoring (SHM) is one of the most
mature approaches for plate-like structures: a sparse network of
piezoelectric transducers excites and receives ultrasonic Lamb waves,
and scattering at damage is detected as a change relative to baseline
records. Deployed long-term, however, the technique faces two coupled
problems. First, the \emph{data deluge}: a monitoring system sampling
continuously over months or years produces gigabytes of records although
the structure is almost always in its nominal state
\cite{yang2025dataset}. Second, \emph{environmental variation}---chiefly
temperature---modifies wave speed, attenuation and transducer coupling
so strongly that temperature-induced signal changes routinely exceed
damage-induced ones \cite{croxford2010,clarke2010}.

The dominant response to the second problem is \emph{temperature
compensation}: baseline signal stretch (BSS) and optimal baseline
selection (OBS) \cite{croxford2010}, physics-based models, and more
recently learned compensators. The dominant response to the first is
\emph{post-hoc compression} of the stored records, e.g.\ SVD/PCA-based
reduction of the long-term data matrix \cite{yang2023compression}. Both
leave the acquisition itself untouched: every record is sampled
uniformly at full rate, whether or not it carries information.

This paper asks a different question: what if the acquisition itself is
\emph{event-based}? Neuromorphic vision has shown that send-on-delta
(SoD) sampling---emitting an event only when the signal changes by a
threshold $\delta$ \cite{miskowicz2006}---preserves task-relevant
information at a small fraction of the data rate. Guided-wave SHM is a
natural candidate: most records are ``nothing happened'', so the
compensated residual is near zero most of the time and an event-based
representation should be extremely sparse. More interestingly,
temperature drift and damage are expected to produce
\emph{qualitatively different event patterns}. Temperature drift is
slow and affects all propagation paths coherently; damage appears
suddenly, persists, and affects only the paths geometrically related to
its location. This suggests that the \emph{temporal and spatial
structure of the event stream}---not the compensated amplitude---can
separate the two, sidestepping explicit temperature compensation.

We make the following contributions:
\begin{itemize}
\item[\textbf{C1}] The first send-on-delta eventization of guided-wave
SHM data, at two levels: within-record level-crossing events of the
compensated residual (microsecond scale), and across-record SoD events
of the damage-index series (day/month scale).
\item[\textbf{C2}] A temperature--damage discrimination mechanism that
uses only event-stream structure (cross-path coherence and burstiness),
without explicit temperature compensation, complementary to the
BSS/OBS paradigm.
\item[\textbf{C3}] A controlled comparison of acquisition-side
eventization against post-hoc compression at equal data rate.
\item[\textbf{C4}] Validation on 4.5 years of real outdoor long-term
data (rain, snow, $-12$ to $+52\,^\circ$C, sensor drift), a regime in
which most guided-wave ML methods are never tested.
\end{itemize}
